%! Author = sriadityavedantam
%! Date = 3/29/26

\section{A variation on Induction}
Now with mathematical induction, also just referenced as induction, we can also show another type called strong or complete induction.
\begin{theorem}{Assume that $n \in \mathbb{Z}^{\geq 0}$ and $P(n)$ is given.
If
    \begin{enumerate}
        \item $P(0)$ is true, and
        \item for all $t \in \mathbb{Z}^{\geq 0}$, if $P(j)$ is true for all $j$ such that $0 \leq j \leq t$, then $P(t+1)$ is also true,
    \end{enumerate}
    then $P(n)$ is true for all $n \in \mathbb{Z}^{\geq 0}$.
}
Let us prove this using ordinary induction.
Let our induction hypothesis be that $P(j)$ is true for all $j$ such that $0 \leq j \leq t$.
Define the set
\[
    S \coloneqq \{n \in \mathbb{Z}^{\geq 0} : P(j) \text{ is true for all } j \text{ such that } 0 \leq j \leq n\}.
\]
For the base case, let $n = 0$.
Then $P(0)$ is true, so $0 \in S$.
Now suppose $k \in S$.
Then $P(j)$ is true for all $j$ such that $0 \leq j \leq k$.
By the hypothesis, this implies $P(k+1)$ is true.
Hence $k+1 \in S$.
Therefore, by induction, $S = \mathbb{Z}^{\geq 0}$.
Thus $P(n)$ is true for all $n \in \mathbb{Z}^{\geq 0}$.
\end{theorem}

Similar to how we used weak or regular induction to prove complete induction, we can do the same in reverse.
In fact, we can prove all of these theorems and definitions using one another.
We can use the well-ordering axiom to prove mathematical induction and use mathematical induction to prove complete induction.
To complete the loop, we can prove well-ordering through complete induction.
On a harder note, we can prove regular induction through complete induction, but it is possible.

\begin{theorem}{Well-Ordering implies Induction.}
Let us define the set
\[
    S \coloneqq \{n \in \mathbb{Z}^{\geq 0} : P(n) \text{ is false}\} \subseteq \mathbb{Z}^{\geq 0}.
\]
Our goal is to show that $S = \varnothing$.
Assume $S \neq \varnothing$.
Then by well-ordering, let $d \in S$ be the smallest element.
Since $P(0)$ is true, we must have $d \neq 0$.
Hence $d \geq 1$, so $d - 1 \in \mathbb{Z}^{\geq 0}$.
Since $d$ is the smallest element of $S$, we must have $d - 1 \notin S$.
Thus $P(d-1)$ is true.
By the induction hypothesis, $P(d-1) \implies P(d)$.
Hence $P(d)$ is true, which contradicts the fact that $d \in S$.
Therefore $S = \varnothing$.
Hence $P(n)$ is true for all $n \in \mathbb{Z}^{\geq 0}$.
\end{theorem}

Now that we have jump-started the proof writing structure in our heads, let us go ahead and start this course with our next topic: Fundamentals of Arithmetic and Divisibility.